\documentclass[a4paper,11pt]{report}
\usepackage{chuan_math_gkt}
\usetikzlibrary{angles,quotes,intersections,patterns,decorations.pathreplacing}
\chuanexamsetup

\begin{document}
\examfontsize

\examsecret{绝密$\bigstar$启用前}

\examtitle{2021年普通高等学校招生全国统一考试（全国乙卷理科）}{数\quad 学}

\noticetitle
\begin{examnotice}
\item 答卷前，考生务必将自己的姓名、准考证号填写在答题卡上。
\item 回答选择题时，选出每小题的答案后，用铅笔把答题卡上对应题目的答案标号涂黑。如需改动，用橡皮擦干净后，再选涂其他答案标号。回答非选择题时，将答案写在答题卡上。写在本试卷上无效。
\item 考试结束后，将本试卷和答题卡一并交回。
\end{examnotice}
\examsection{一、选择题：本题共12小题，每小题5分，共60分。在每小题给出的四个选项中，只有一项是符合题目要求的。}
\begin{examquestions}

\item 
设$2(z+\overline{z})+3(z-\overline{z})=4+6i$，则$z=$

\fourchoices{$1-2i$}{$1+2i$}{$1+i$}{$1-i$}

\item 
已知集合$S=\{s|s=2n+1,n\ensuremath{\in} Z|\}$，$T=\{t|t=4n+1,n\ensuremath{\in} Z|\}$，则$S\cap T=$

\fourchoices{$\varnothing$}{$S$}{$T$}{$Z$}

\item 
已知命题$p:∃x\ensuremath{\in} R,sinx<1$﹔命题$q:∀x\ensuremath{\in} R$﹐$e^{|x|}\ensuremath{\geq}1$，则下列命题中为真命题的是

\fourchoices{$p∧q$}{$¬p∧q$}{$p∧¬q$}{$¬(p∨q)$}

\item 
设函数$f(x)=\dfrac{1-x}{1+x}$，则下列函数中为奇函数的是

\fourchoices{$f(x-1)-1$}{$f(x-1)+1$}{$f(x+1)-1$}{$f(x+1)+1$}

\item 
在正方体$ABCD-A_{1}B_{1}C_{1}D_{1}$中，\emph{P}为$B_{1}D_{1}$的中点，则直线$PB$与$AD_{1}$所成的角为

\fourchoices{$\dfrac{\pi}{2}$}{$\dfrac{\pi}{3}$}{$\dfrac{\pi}{4}$}{$\dfrac{\pi}{6}$}

\item 
将5名北京冬奥会志愿者分配到花样滑冰、短道速滑、冰球和冰壶4个项目进行培训，每名志愿者只分配到1个项目，每个项目至少分配1名志愿者，则不同的分配方案共有

\fourchoices{60种}{120种}{240种}{480种}

\item 
把函数$y=f(x)$图像上所有点的横坐标缩短到原来的$\dfrac{1}{2}$倍，纵坐标不变，再把所得曲线向右平移$\dfrac{\pi}{3}$个单位长度，得到函数$y=\sin(x-\dfrac{\pi}{4})$的图像，则$f(x)=$

\onechoices{$\sin(\dfrac{x}{2}-\dfrac{7\pi}{12})$}{$\sin(\dfrac{x}{2}+\dfrac{\pi}{12})$}{$\sin(2x-\dfrac{7\pi}{12})$}{$\sin(2x+\dfrac{\pi}{12})$}

\item 
在区间$(0,1)$与$(1,2)$中各随机取1个数，则两数之和大于$\dfrac{7}{4}$的概率为

\fourchoices{$\dfrac{7}{9}$}{$\dfrac{23}{32}$}{$\dfrac{9}{32}$}{$\dfrac{2}{9}$}

\item 
魏晋时刘徽撰写的《海岛算经》是有关测量的数学著作，其中第一题是测海岛的高．如图，点$E$，$H$，$G$在水平线$AC$上，$DE$和$FG$是两个垂直于水平面且等高的测量标杆的高度，称为``表高''，$EG$称为``表距''，$GC$和$EH$都称为``表目距''，$GC$与$EH$的差称为``表目距的差''则海岛的高$AB=$

$\dfrac{\text{表高}\times\text{表距}}{\text{表目距的差}}+\text{表高}$

\onechoices{$\dfrac{\text{表高}\times\text{表距}}{\text{表目距的差}}+\text{表高}$}{$\dfrac{\text{表高}\times\text{表距}}{\text{表目距的差}}-\text{表高}$}{$\dfrac{\text{表高}\times\text{表距}}{\text{表目距的差}}+\text{表距}$}{$\dfrac{\text{表高}\times\text{表距}}{\text{表目距的差}}-\text{表距}$}

\item 
设$a\ne0$，若$a$为函数$f(x)=a(x-a)^{2}(x-b)$的极大值点，则

\fourchoices{$a<b$}{$a>b$}{$ab<a^{2}$}{$ab>a^{2}$}

\item 
设$B$是椭圆$C:\dfrac{x^{2}}{a^{2}}+\dfrac{y^{2}}{b^{2}}=1(a>b>0)$$F$上顶点，若$C$上的任意一点$P$都满足$|PB|\ensuremath{\leq}2b$，则$C$的离心率的取值范围是

\onechoices{$(\dfrac{\sqrt{2}}{2},1)$}{$(\dfrac{1}{2},1)$}{$(0,\dfrac{\sqrt{2}}{2})$}{$(0,\dfrac{1}{2})$}

\item 
设$a=2ln1.01$，$b=ln1.02$，$c=\sqrt{1.04}-1$．则

\fourchoices{$a<b<c$}{$b<c<a$}{$b<a<c$}{$c<a<b$}

\item 
已知双曲线$C:\dfrac{x^{2}}{m}-y^{2}=1(m>0)$的一条渐近线为$\sqrt{3}x+my=0$，则\emph{C}的焦距为\_\_\_\_\_\_\_\_\_．

\item 
已知向量$\vec{a}=(1,3),\vec{b}=(3,4)$，若$(\vec{a}-\lambda\vec{b})\ensuremath{\perp}\vec{b}$，则$\lambda=$\blank．

\item 
记$\ensuremath{\triangle}ABC$的内角\emph{A}，\emph{B}，\emph{C}的对边分别为\emph{a}，\emph{b}，\emph{c}，面积为$\sqrt{3}$，$B=60^\circ$，$a^{2}+c^{2}=3ac$，则$b=$\_\_\_\_\_\_\_\_．

\item 
以图①为正视图，在图②③④⑤中选两个分别作为侧视图和俯视图，组成某个三棱锥的三视图，则所选侧视图和俯视图的编号依次为\_\_\_\_\_\_\_\_\_（写出符合要求的一组答案即可）．

$\overline{x}$

\end{examquestions}
\examsection{三、解答题：共70分。解答应写出文字说明、证明过程或演算步骤。第17题～第21题为必考题，每个试题考生都必须作答。第22、23题为选考题，考生根据要求作答。}
\begin{examanswers}[start=17]


\item（12分）

某厂研制了一种生产高精产品的设备，为检验新设备生产产品的某项指标有无提高，用一台旧设备和一台新设备各生产了10件产品，得到各件产品该项指标数据如下：

\begin{tightcenter}
\begin{tabular}{|c|c|c|c|c|c|c|c|c|c|c|}
\hline
旧设备 & 9.8 & 10.3 & 10.0 & 10.2 & 9.9 & 9.8 & 10.0 & 10.1 & 10.2 &
97 \\ \hline
新设备 & 10.1 & 10.4 & 10.1 & 10.0 & 10.1 & 10.3 &
106
& 10.5 & 10.4 & 10.5 \\ \hline
\end{tabular}
\end{tightcenter}

旧设备和新设备生产产品的该项指标的样本平均数分别记为$\overline{x}$和$\overline{y}$，样本方差分别记为$s_{1}^{2}$和$s_{2}^{2}$．

（1）求$\overline{x}$，$\overline{y}$，$s_{1}^{2}$，$s_{2}^{2}$；

（2）判断新设备生产产品的该项指标的均值较旧设备是否有显著提高（如果$\overline{y}-\overline{x}\ensuremath{\geq}2\sqrt{\dfrac{s_{1}^{2}+s_{2}^{2}}{10}}$，则认为新设备生产产品的该项指标的均值较旧设备有显著提高，否则不认为有显著提高）．

\item（12分）

如图，四棱锥$P-ABCD$$F$底面是矩形，$PD\ensuremath{\perp}$底面$ABCD$，$PD=DC=1$，$M$为$BC$的中点，且$PB\ensuremath{\perp} AM$．

$BC$

（1）求$BC$；

（2）求二面角$A-PM-B$的正弦值．

\item（12分）

记$S_{n}$为数列$\{a_{n}\}$的前\emph{n}项和，$b_{n}$为数列$\{S_{n}\}$的前\emph{n}项积，已知$\dfrac{2}{S_{n}}+\dfrac{1}{b_{n}}=2$．

（1）证明：数列$\{b_{n}\}$是等差数列;

（2）求$\{a_{n}\}$的通项公式．

\item（12分）

设函数$f(x)=\ln(a-x)$，已知$x=0$是函数$y=xf(x)$的极值点．

（1）求\emph{a}；

（2）设函数$g(x)=\dfrac{x+f(x)}{xf(x)}$．证明:$g(x)<1$．

\item（12分）

已知抛物线$C:x^{2}=2py(p>0)$$F$焦点为$F$，且$F$与圆$M:x^{2}+(y+4)^{2}=1$上点的距离的最小值为$4$．

（1）求$p$；

（2）若点$P$在$M$上，$PA,PB$是$C$的两条切线，$A,B$是切点，求$\ensuremath{\triangle}PAB$面积的最大值．


{{[}选修4-4：坐标系与参数方程{]}（10分）}

\item（10分）

在直角坐标系$xOy$中，$\odot C$的圆心为$C(2,1)$，半径为1．

（1）写出$\odot C$的一个参数方程;

（2）过点$F(4,1)$作$\odot C$的两条切线．以坐标原点为极点，\emph{x}轴正半轴为极轴建立极坐标系，求这两条切线的极坐标方程．

{{[}选修4-5：不等式选讲{]}（10分）}

\item（10分）

已知函数$f(x)=|x-a|+|x+3|$．

（1）当$a=1$时，求不等式$f(x)\ensuremath{\geq}6$的解集;

（2）若$f(x)>-a$，求\emph{a}的取值范围．
\end{examanswers}

\end{document}
